\subsection{Arquitectura en cascada}

El controlador clásico usa un lazo externo de posición para producir una fuerza
deseada en mundo ENU y un lazo interno para aproximar la actitud, el empuje y los
momentos necesarios. Esta división se conserva después en el controlador
neuronal para mantener una base común de estabilización angular.

\subsection{Ecuaciones, ganancias y límites}

Se desarrollarán los errores de posición, velocidad y actitud, así como las
ganancias y saturaciones efectivas. Aunque se utiliza la denominación PID en la
comparación, la implementación actual corresponde a una estructura PD en cascada
con términos de referencia y compensación gravitatoria.

\subsection{Alternativas de sintonización}

Se revisarán ajuste manual, Ziegler--Nichols, métodos basados en modelo y
búsquedas numéricas. La elección final se justificará por la necesidad de
evaluar candidatos directamente sobre escenarios simulados reproducibles y con
restricciones de seguridad.

\subsection{Sintonización y congelación de controladores}

El tuneo progresivo determinista utiliza únicamente escenarios de entrenamiento,
aplica filtros duros y registra semilla, presupuesto y criterio de selección.
Los PIDs aceptados se congelan por familia antes de generar la evidencia de
baseline y de transferencia.

\subsection{Baseline y transferencia cruzada}

La comparación distingue PIDs especializados, PIDs aplicados a otras familias y
un baseline representativo. Esta matriz permite separar rendimiento ajustado de
capacidad de transferencia y proporciona una referencia adecuada para valorar
las políticas neuronales.
